\documentclass[../../../include/open-logic-section]{subfiles}

\begin{document}

\olfileid{sth}{story}{rus}
\olsection{రసెల్ వైరుధ్యం (మళ్లీ)}

\olref[sfr][][]{part}లో మనం
అనౌపచారిక సమితి సిద్ధాంతంతో
పనిచేశాం. అయితే \emph{అతి}
అమాయక భావన ప్రకారం సమితులు
విధేయాలు వర్తించే వస్తువుల
సమితులే. ఆ అమాయక భావన
కింది సూత్రాన్ని కోరుతుంది:

\begin{defish}
  \emph{నిర్బంధరహిత ధర్మసంగ్రహం.}
  ఏ సూత్రం $\phi$కైనా
  $\Setabs{x}{\phi(x)}$ ఉంటుంది.
\end{defish}

ఈ సూత్రం ఆకర్షణీయమైనప్పటికీ,
ఇది వైరుధ్యపూరితమని
నిరూపించవచ్చు. దీనిని
\olref[sfr][set][rus]{sec}లో
చూశాం. కానీ ఫలితం ఎంతో
ముఖ్యమైనదీ, నేరుగా అర్థమయ్యేదీ
కాబట్టి, అచ్చంగా మళ్లీ
చెప్పడం సముచితం.

\begin{thm}[రసెల్ వైరుధ్యం]
$R = \Setabs{x}{x \notin x}$
అనే సమితి లేదు.
\end{thm}

\begin{proof}
$R = \Setabs{x}{x \notin x}$
ఉందనుకుంటే, $R \in R$
అప్పుడే $R \notin R$;
ఇది వైరుధ్యం.
\end{proof}

రసెల్ ఈ ఫలితాన్ని జూన్ 1901లో
కనుగొన్నారు. (అయితే మనం ఇప్పుడు
ఇచ్చిన రూపంలోనే ఆయన వైరుధ్యాన్ని
ప్రతిపాదించలేదు; ఎందుకంటే ఆయన
\emph{గ్రుండ్‌గెజెట్సె}లో వివరించిన
ఫ్రేగె సమితి సిద్ధాంతాన్ని
పరిశీలిస్తున్నారు. దీనికి
\olref[blv]{sec}లో తిరిగి వస్తాం.)
ఫ్రేగె వ్యవస్థలోని వైరుధ్యాన్ని
వివరిస్తూ రసెల్ 1902 జూన్ 16న
ఆయనకు లేఖ రాశారు. ఆ ఉత్తర
ప్రత్యుత్తరాలకూ కొంత నేపథ్యానికీ
\citet[pp.~124--8]{Heijenoort1967}
చూడండి.

ఈ రెండు పంక్తుల నిరూపణ
\emph{శుద్ధ తర్కం} నుంచి
వస్తుందని నొక్కి చెప్పాలి.
నిజమే, $\Setabs{x}{x \notin x}$
అనే సంకేతాన్ని వాడటంలో
సమితుల సమానత్వం అనే
(తార్కికం కాని?)\ స్వీకృతాన్ని
అంతర్లీనంగా వాడాం;
ఎందుకంటే $\Setabs{x}{\phi(x)}$
ఉంటే, అది $\phi$ను
సంతృప్తిపరిచే వస్తువుల
\emph{ఏకైక} సమితి
(సమితుల సమానత్వం ప్రకారం).
అయితే ఫలితాన్ని ఇలా
చెబితే ఆ స్వీకృత సూచన
కూడా అవసరం లేదు:
\emph{తమలో తాము మూలకాలుగా
లేని సమితులే ఖచ్చితంగా
మూలకాలుగా ఉన్న సమితి లేదు}.
నిజానికి దీనికి సమితులతో
ప్రత్యేక సంబంధమేమీ లేదు.
రసెల్ స్వయంగా గమనించినట్లు,
అచ్చం ఇలాంటి వాదనే
\emph{తమను తాము క్షౌరం
చేసుకోని పురుషులనే ఖచ్చితంగా
క్షౌరం చేసే పురుషుడు లేడు}
అనే ముగింపునిస్తుంది.
లేదా \emph{తమను తాము
వాసన చూడని పగ్ కుక్కలనే
ఖచ్చితంగా వాసన చూసే
పగ్ కుక్క లేదు}.
ఇలాగే మరెన్నో. సంకేతరూపంలో
ఫలితం ఆకృతి కేవలం ఇది:
\[
\lnot \exists x \forall z(Rzx \liff \lnot R zz).
\]
ఇది మొదటిస్థాయి తర్కంలోని
ఒక సిద్ధాంతం (లేదా సిద్ధాంత
పథకం) మాత్రమే. అందువల్ల
మన సమితి సిద్ధాంతాన్ని
కొద్దిగా సవరించడం ద్వారా
మాత్రమే రసెల్ వైరుధ్యాన్ని
నివారించలేం; అది సమితి
సిద్ధాంతానికి \emph{చేరకముందే}
వస్తుంది. మనం (సాంప్రదాయ)
మొదటిస్థాయి తర్కాన్ని వాడాలనుకుంటే,
$R = \Setabs{x}{x\notin x}$
అనే సమితి లేదని
\emph{అంగీకరించాల్సిందే}.

సారాంశం ఇది. వైరుధ్యాన్ని
\emph{నివారిస్తూ} నిర్బంధరహిత
ధర్మసంగ్రహాన్ని అంగీకరించాలంటే,
\emph{సమితి సిద్ధాంతం}ను
మాత్రమే స్వల్పంగా సవరించడం
సరిపోదు. దాని బదులు
\emph{తర్కం}నే మూలంగా
మార్చాలి.

సాంప్రదాయేతర తర్కాలతో
సమితి సిద్ధాంతాలను
ప్రతిపాదించారు. కానీ అవి,
కనీసం చెప్పాలంటే,
ప్రామాణికమైనవి కావు.
రసెల్ వైరుధ్యాన్ని నేరుగా
అస్తిత్వరాహిత్య నిరూపణగా
తీసుకొని, దానితో ఎలా
వ్యవహరించాలో నేర్చుకోవడమే
ప్రామాణిక పద్ధతి. మనమూ
అదే పద్ధతిని అనుసరిస్తాం.

\end{document}
